%%%%%%%%%%  scan idx 229  |  folio 218  %%%%%%%%%%

\begin{center}
SGA 7

\vspace{1em}
\underline{Expos\'e VIII}

\vspace{1em}
\underline{\underline{COMPLEMENTOS SOBRE LAS BIEXTENSIONES. PROPIEDADES GENERALES}}

\underline{\underline{DE LAS BIEXTENSIONES DE ESQUEMAS EN GRUPOS}}

\vspace{1em}
\underline{por A. Grothendieck}
\end{center}

\subsection*{Sumario}

\begin{enumerate}[label=\arabic*.]
\setcounter{enumi}{-1}
\item Introducción
\item Diversos casos particulares sobre un topos cualquiera
\item Acoplamientos definidos por una biextensión
\item Biextensiones de esquemas en grupos $(P,Q)$ por $\underline G_m$: generalidades
\item Extensiones y biextensiones de esquemas en grupos lisos y conexos sobre un cuerpo
\item Extensiones y biextensiones por esquemas en grupos constantes truncados sin torsión
\item Extensiones y biextensiones por el modelo de Néron de $\underline G_m$
\item Prolongaciones canónicas de extensiones y biextensiones por $\underline G_m$
\end{enumerate}

0. \underline{Introducción}

0.1. Continuamos aquí el estudio general de las biextensiones iniciado
en el exposé precedente. Damos en primer lugar algunos complementos en
el caso de un topos de base cualquiera, desarrollando en particular,
con el detalle que merece, el formalismo de los acoplamientos asociados
a las biextensiones (sección 2), cuya importancia en la teoría de los
esquemas abelianos es bien conocida. Abordamos después el estudio de
las biextensiones de esquemas en grupos generales; el caso que más nos
interesa es el de las biextensiones por el grupo multiplicativo
$\underline G_m$. No obstante, como etapa técnica intermedia, nos veremos
obligados a estudiar también las biextensiones por otros tipos de
esquemas en grupos (secciones 5 y 6).

%%%%%%%%%%  scan idx 230  |  folio 219  %%%%%%%%%%

El resultado más interesante para lo que sigue se da en 7.1. Allí se
establecen condiciones bajo las cuales, para $P$ y $Q$ lisos sobre un
rasgo de punto genérico $\eta$, una biextensión de
$(P_\eta,Q_\eta)$ por $\underline G_{m\eta}$ puede prolongarse a una
biextensión de $(P,Q)$ por $\underline G_m$, determinada salvo
isomorfismo único. Este resultado será muy cómodo en el exposé
siguiente, dedicado a la teoría del modelo de Néron de las variedades
abelianas. Es probable que allí pudiera llegarse al resultado con
menos aparato, sin el imponente pisapapeles de los exposés VII y VIII
(que ya no se utilizarán en el resto del Seminario); con todo, nos
parece probable que los desarrollos dados en estos exposés estén
destinados a ser útiles también en otros lugares.

0.2. Cuando trabajemos con esquemas en grupos sobre un esquema $S$, se
sobreentenderá que trabajamos con el topos fppf de $S$. Este es el
topos asociado al $\underline U$-sitio de los esquemas localmente de
presentación finita sobre $S$ (y pertenecientes a $\underline U$),
dotado de la topología fppf (SGA 3 IV 6.3). El lector comprobará sin
dificultad que la mayoría de los enunciados que daremos serían
igualmente válidos (y esencialmente equivalentes a los enunciados
fppf) para otras topologías sobre $S$ cuyo sitio subyacente contenga
al menos los esquemas localmente de presentación finita sobre $S$,
como las topologías de Zariski, étale y fpqc. Una de las razones es que,
por la teoría del descenso, la categoría de las extensiones de un
esquema en grupos $P$, localmente de presentación finita sobre $S$,
por el esquema en grupos $\underline G_m$ no depende, salvo
equivalencia, de la topología elegida; lo mismo vale para la noción de
biextensión.

La razón técnica para preferir la topología fppf (o una topología más
fina, como la fpqc) a topologías menos finas como las de Zariski y
étale es que, para todo entero $n>0$, el homomorfismo $n\cdot\id$ de
$\underline G_m$ en sí mismo es un \underline{epimorfismo fppf}
(pues es fielmente plano y de presentación finita), pero en general no
es un epimorfismo de Zariski ni étale. Obsérvese que es un epimorfismo
étale si y solo si $n$ es primo con las características residuales de
$S$, y un epimorfismo de Zariski si y solo si $S$ es vacío.

%%%%%%%%%%  scan idx 231  |  folio 220  %%%%%%%%%%

0.3. Señalemos finalmente que la mayoría de los enunciados relativos a
extensiones o biextensiones por $\underline G_m$ seguirían siendo
válidos si se reemplazara este grupo por un \underline{toro} $G$.
Como tal $G$ es localmente isomorfo (para la topología fppf, e incluso
para la topología étale) a un grupo $\underline G_m^r$, se trata de una
generalización esencialmente trivial.
